% !PRINTED 732
\weberpage{732}
%<<CONT
behave in a certain sense like a fundamental system, which is expressed
more definitely by the following theorem:

2. \gesperrt{If a normal unit of the field $H_m$
of the form}
\[
\mathfrak{E}(r) = \tau_0^{e_0}\tau_1^{e_1}\ldots\tau_{\nu'-1}^{e_{\nu'-1}},
\]
\gesperrt{in which the exponents $e_0, e_1, \ldots, e_{\nu'-1}$ are integral rational
 numbers, has one and the same sign as all its conjugate
values, then all of the $e_0, e_1, \ldots, e_{\nu'-1}$ must be even numbers,
and $\mathfrak{E}(r)$ is the square of a normal unit.}

We can again carry out the proof by complete induction, but wish first
to transform the expression for $\tau_\beta$ somewhat.

We obtain the conjugate values to the unit $\mathfrak{E}(r)$
in the form
\[
\mathfrak{E}(r_\beta) = \tau_\beta^{e_0}\tau_{\beta+1}^{e_1}\ldots\tau_{\beta+\nu'-1}^{e_{\nu'-1}},
\]
and these expressions are therefore to have the same sign for all
values of $\beta$.

By §.\,199, (5) we have
\[
\tau_\beta = \mathrm{tg}\,\frac{5^\beta\pi}{m} = \frac{1}{2}\frac{\sin\left(\dfrac{2\pi}{m}5^\beta\right)}{\left[\cos\left(\dfrac{5^\beta\pi}{m}\right)\right]^2}.
\]

In this we put, for abbreviation,
\weq{14}{\sigma_\beta = \sin\left(\frac{2\pi}{m}5^\beta\right),}
so that $\sigma_\beta$ always has the same sign as $\tau_\beta$.

From this we form the product
\weq{15}{S_\beta = \sigma_\beta^{e_0}\sigma_{\beta+1}^{e_1}\ldots\sigma_{\beta+\nu'-1}^{e_{\nu'-1}},}
and our assumption is therefore this, that $S_\beta$ has the same sign
for all values of $\beta$.

It is to be shown that the integers
$e_0, e_1, \ldots, e_{\nu'-1}$ must all be even. The theorem is
correct for $m=8$; for in this case we have
\[
\sigma_0 = \sin\frac{2\pi}{8}, \sigma_1 = \sin\frac{10\pi}{8} = -\sin\frac{6\pi}{8},
\]
%CONT>>
